\section{Limitations}
\label{sec:limitations}
\begin{itemize}
\item \textbf{No human labels.} Every rate in this draft uses labels proposed by a pre-specified rule. The rule combines construction intent, model-written reference checks and the repository's own tests. Where these signals disagree, the candidate is excluded as ambiguous (Table~\ref{tab:labels}). An independent model check disagrees with several proposals, mostly on edge cases (exotic Unicode input, a read-only file mid-session). Whether such cases violate a requirement is a human judgement. The paper's central claim needs human labels and is not supported by this pilot.
\item \textbf{Controlled population.} The four sessions are real agent runs, but the study wrote their user turns and designed each to contain a correction. The correction reached the agent through a messaging channel that adds harness framing. Natural correction frequency and content are not represented.
\item \textbf{The anchor negative is not always a negative.} In some sessions the agent's pre-correction state already satisfied most or all of the correction. That case is informative, but it means that ``the state the user rejected'' is not automatically an invalid candidate. Some later corrections restate what a careful agent would have done anyway.
\item \textbf{The information lost depends on the contract proposer.} For one task the initial-only contract already contained, as \emph{inferred} requirements, most of what the correction later stated (Table~\ref{tab:coverage}). The initial-only arm therefore measures what one proposer recovers from the first request, not a fixed property of the session.
\item \textbf{One model family.} The sessions, contracts, verifiers, attacks, alternatives, near misses, reference checks, label proposals, construction checks and both requirement mappers come from models of one family. This can correlate errors across conditions and with the labels. The high agreement between the two mappers is same-model agreement.
\item \textbf{Near misses are built from the reference.} The generators started from the reference solution and introduced one plausible violation. Such near misses may be easier or harder to catch than the errors agents make naturally.
\item \textbf{Four task families.} Intervals over four families describe variability; they are not evidence of an effect, and no comparison is tested for significance.
\item \textbf{Harness defects during the pilot.} Two resource-isolation defects were found after the first evaluation run (\S\ref{sec:method}). Both are fixed, and all reported results were re-run under the corrected harness. The development audits ran under the earlier harness versions.
\item \textbf{Asymmetric reference-check review.} Only the last task's reference checks went through an independent review-and-fix loop before freezing. The others were frozen, and used for labels, before that loop existed.
\item \textbf{Exploit surface.} Within one sandbox call, candidate code controls the value it reports. The repository-test baseline imports candidate code into the reporting process, and a deliberately adversarial candidate can forge its report (\S\ref{sec:repro}).
\item \textbf{Contamination.} The source repositories and the released tasks are public and may enter future training data.
\end{itemize}
